\subsection{The normalized actual conjugator belongs to the simple socle}

\begin{statement}{D23.F32}{Proved}{Eliminating the remaining binary block coordinates}
\end{statement}
Let $h=(1,f)$ be an actual full-GT element with
\[
                         \rho(h)=0\pmod3,\qquad a_2(h)=0.
\]
For either value of $a_3(h)$, the actual degree-$1344$
Hurwitz conjugator $V_h$ of D23.P18 belongs to
\[
                 S^{12}=\operatorname{P\Omega}_6^-(3)^{12}
                           \subset\Gamma.
\]
In particular this applies to the normalized hypothetical
mixed state $(\rho,a_3,a_2)=(0,1,0)$ of D23.H2. The statement
is a necessary condition, not a contradiction for that state.

\begin{proof}
First determine the exact signed solvable values, retaining
the central ambiguity until it is resolved. By D23.F31, the
signed $24$-point quotient is marked-isomorphic to the
matrices FP.36.2 modulo $\langle\omega I\rangle$.
Let $\overline F$ and $\overline R$ be the projective Burau
values of $f(X_1,X_2)$ and $f(X_3,X_4)$.
The supplied FP.37 uses exactly these positive half twists,
$X_1=A^2$, $X_2=B^2$, $X_3=C^2$, and
$X_4=(ABA)^2$. Its actual-full-GT consequences are
\[
 a_f=\rho\pmod3,\qquad \nu_4=a_2,
\]
and its exact finite formulas give
\[
 \overline F=F^0_{\rho,a_2}Z^c,
 \qquad
 \overline R=(e_0t_0^2)^{\nu_4}p_0^{a_f},
 \qquad c\in\mathbb F_2.
\]
Here $Z=(ABA)^4$ in the first-pair image, and the literal
table of FP.37 has $F^0_{0,0}=I$. For the present normalized
state, $a_f=\nu_4=0$, hence $\overline R=I$ and
$\overline F=Z^c$. The same source's exact pentagon formula
(FP.37.6) is
\[
                     I=\mathcal P(f)=w_0^c
\]
in this case. Its displayed $w_0$ is a nonidentity involution
(its matrix has nonzero off-diagonal entries), so $c=0$.
Thus both projective values are the identity. The determinant
cube character of D23.F31 is one on every commutator
evaluation and is minus one on the remaining central scalar
$-I$. The signed values are consequently
\[
                             F_{24}=R_{24}=I.
\]
This uses the actual pentagon in the supplied Burau quotient;
the conclusions are not imposed on arbitrary gentle residues.

It follows from D23.P18 that the induced conjugation of the
signed quotient fixes each of its three marked generators.
In particular the induced permutation of the twelve blocks
centralizes $Q_{12}$. That transitive group's point-stabilizer
orbits have sizes $1,2,9$, as checked in D23.F30. Its
permutation centralizer is therefore trivial: the image of
one point under a commuting permutation must be a fixed
point of its stabilizer, and there is exactly one. Thus
$V_h$ fixes all twelve blocks.

Use the group-transversal identifications of D23.F30, and
write $V_i$ for the local permutation on block $i$.
D23.H3 says that, for either normalized cubic value, the
four local maps on blocks $7,8,10,11$ belong to $S$.
It suffices below to use just one such block.
For any generator $g$ taking block $i$ to block $j$, let
$L(g,i)=T_j^{-1}gT_i\in L$ be its local transition.
The actual covariance gives
\[
            V_jL(g,i)=L(\Phi_h(g),i)V_i,
\]
where both transition maps belong to $L$ because
$g,\Phi_h(g)\in\Gamma$ and the blocks are fixed by $V_h$.
Starting at a block with $V_i\in S\subset L$ and traversing
the transitive block action proves $V_j\in L$ for every $j$.
This propagation does not assume in advance that $V_h$
belongs to the finite braid image.

Let $v_i\in L/S=\mathbb F_2$ be the twelve local cosets.
Passing the last displayed equation to $L/S$, the equality
of the signed generator and its transformed signed value
gives
\[
                          v_{q_g(i)}=v_i
                \qquad(g=A,B,C).
\]
Indeed the two transition characters are identical because
$F_{24}=R_{24}=I$ and $\Phi_h(A)=A$ in that quotient.
Transitivity of $Q_{12}$ makes all $v_i$ equal. One is zero
by D23.H3, so every $v_i$ is zero. Thus every local $V_i$
belongs to $S$. D23.F30 identifies the independent product
of these twelve local groups with an actual subgroup
$S^{12}\subset\Gamma$. Therefore $V_h\in S^{12}$, as
claimed.
\end{proof}

\paragraph{Exact dependency scope.}
The only supplied arithmetic normalization used to reach this
case is D23.H2. The finite Burau input is supplied FP.37,
especially its explicit first-pair table and equations
(FP.37.4)--(FP.37.6), in the original positive marking.
D23.F31 removes the possible scalar lift ambiguity. The
remaining proof uses the actual covariance D23.P18, the
checked block structure D23.F30, and the exhaustive local
lift theorem D23.H3. It does not infer an actual GT element
from finite solutions and makes no half-Kummer comparison.
